\documentclass[preprint,12pt]{elsarticle}
\usepackage{amsmath,amssymb,array,booktabs,graphicx,lineno,tabularx,xurl}
\usepackage{hyperref}
\hypersetup{colorlinks=true,linkcolor=blue,citecolor=blue,urlcolor=blue}
\biboptions{sort&compress}
\newcolumntype{Y}{>{\raggedright\arraybackslash}X}
\setlength{\emergencystretch}{3em}
\journal{Journal of Safety Science and Resilience}
\begin{document}
\begin{frontmatter}
\title{Provenance-Preserving Evidence-Gated Knowledge-Graph Augmentation for Auditable Entity--Relation Extraction}
\author{Anonymous Authors}
\begin{abstract}
\textbf{Introduction:} Graph-augmented extraction can retrieve useful
associations, but a retrieved association is not necessarily a fact stated in
the current sentence. Fluent generators also make it difficult to trace a
predicted relation to its source span and to audit why it was accepted. This
gap is consequential when extracted assertions are used in safety-related
knowledge systems.\\
\textbf{Methods:} We develop a provenance-preserving, evidence-gated
framework that separates graph-informed candidate generation from output
acceptance. A graph derived only from training annotations supplies bounded
concept and relation hints to QLoRA-adapted Qwen3-4B extractors. Deterministic
span reconstruction, entity acceptance, and relation verification then retain
only source-linked, type-compatible assertions. We evaluate the framework on
ADE and CoNLL04 using a shared strict character-span protocol and report the
underlying dataset distributions.\\
\textbf{Results:} The completed test comparisons show that the provenance-gated
system improves relation F1 over its matched source-only generator on both
benchmarks and records explicit acceptance decisions. The relation improvement
is retained for the reported CoNLL04 sensitivity analysis after exact
train--test text overlaps are excluded.\\
\textbf{Discussion:} The results support controlled graph augmentation as an
auditable extraction layer for human review. The gates establish structural
admissibility and source linkage, not semantic truth, clinical causality, or
universal accuracy dominance; the strongest supervised baseline remains
competitive in overall extraction F1.
\end{abstract}
\begin{keyword}
entity and relation extraction \sep knowledge graph \sep provenance
\sep evidence grounding \sep adverse drug effects \sep selective prediction
\end{keyword}
\end{frontmatter}
\linenumbers

\section{Introduction}
\label{sec:introduction}
Structured information extraction connects unstructured text to searchable
entities, typed relations, and downstream knowledge systems. In applications
related to safety, an extracted association must also be inspectable: users
need to know which mention was identified, which sentence supported the
relation, and which transformations produced the final assertion. Drug-related
adverse-effect extraction illustrates this requirement directly. A model that
recognizes a drug and an adverse effect but invents an association between them
can create a misleading record even when its output is syntactically valid.
General-domain extraction provides a complementary test of the same technical
issue across a broader set of entity and relation labels.

Neural entity and relation models provide effective approaches to the
recognition problem~\cite{li2022ner,eberts2020spert}. Generative models offer a
flexible alternative in which a common structured output format can be reused
across label inventories. However, fluent generation does not guarantee that
the predicted objects are supported by their input~\cite{ji2023hallucination}.
This distinction matters when extracted triples are inserted into a graph:
errors become available to later retrieval or reasoning stages, potentially
propagating beyond the sentence that generated them.

Knowledge graphs can contribute typed domain priors~\cite{hogan2021kg}, and
retrieval augmentation can expose those priors without requiring all knowledge
to be stored in model parameters~\cite{lewis2020rag}. Their use introduces a
specific authority problem. A relation observed in a training sentence may
suggest what to search for, but it does not establish that the same relation
holds in a new sentence. A training graph may also return an example's own
annotation or reveal content shared across data splits. The useful question is
therefore not only whether graph context increases recall, but whether its
eligibility and influence can be controlled and audited.

We address this problem with a staged architecture. Separate source-only,
exact-anchor, and high-recall graph-conditioned generators establish a matched
comparison of input conditions. A deterministic relation verifier checks
structural and source-evidence requirements. A complementary entity gate uses
agreement, source-matched anchors, or verified relation endpoints to accept a
subset of the graph-conditioned candidates. The final system combines the
accepted entities with verified relations whose endpoints survive. The
architecture uses existing representation and adaptation methods; its
contribution is their governed composition, not a new backbone or a new
optimization algorithm.

This paper concentrates on ADE~\cite{gurulingappa2012ade} and
CoNLL04~\cite{roth2004conll}. ADE supplies a drug-safety-related extraction
task with two entity labels and one relation label. CoNLL04 supplies a
general-domain control with four entity labels and five directed relation
labels. Both use English sentence-level inputs, making it possible to examine
the extraction and verification mechanisms without importing a long-document
windowing problem. They are trained and evaluated separately; this is not a
cross-dataset zero-shot transfer experiment.

The contributions are threefold:
\begin{itemize}
\item A provenance-aware graph-context policy that uses only training
annotations, excludes candidates supported only by the current example, and
clears graph hints for exact train/evaluation text overlaps.
\item A deterministic composition of source-span reconstruction, entity
selection, and relation verification that retains evidence-linked assertions
and records the reasons behind output decisions.
\item A reproducible two-benchmark analysis that combines split and label
statistics, uniform strict-span scoring, validation ablations, completed test
comparisons, and an exact-overlap sensitivity analysis.
\end{itemize}
The intended contribution is an auditable extraction layer for analyst-facing
knowledge systems. It does not convert benchmark relation labels into clinical
causality, establish calibrated risk probabilities, or validate autonomous
safety decisions.

\section{Related work}
\label{sec:related}
\subsection{Joint entity and relation extraction}
Entity recognition must resolve both mention boundaries and types; relation
extraction additionally depends on recovering the correct directed
endpoints~\cite{li2022ner}. Global consistency constraints have a long history
in this setting, including the formulation associated with the CoNLL04
benchmark~\cite{roth2004conll}. SpERT enumerates spans and jointly models
entities and relations using a pretrained transformer~\cite{eberts2020spert}.
It is a particularly relevant comparison because it performs the same joint
task without requiring a generative JSON interface. Our experiments retrain
SpERT from a clean base encoder on the same training partitions rather than
using a released task checkpoint that may already include development data.

GLiNER and GLiREL provide label-conditioned alternatives for entity recognition
and relation classification~\cite{zaratiana2024gliner,boylan2025glirel}.
These models differ from locally fine-tuned generators in supervision,
pretraining, and output construction. We therefore report the implemented
train-calibrated pipeline as a named external baseline, not as an exhaustive
assessment of either model family. A same-backbone zero-shot generator further
separates the benefit of local adaptation from that of graph conditioning.

\subsection{Graph context and constrained generation}
Knowledge graphs represent typed entities and relations explicitly, while
retrieval augmentation uses selected records as model
context~\cite{hogan2021kg,lewis2020rag}. GenIE constrains generative information
extraction with a structured output space and entity/relation
catalogues~\cite{josifoski2022genie}. Catalogue validity and source support are
different requirements: an admissible entity or relation may still not be
asserted in the current input. Our graph context is consequently treated as a
candidate prior whose eligibility depends on training provenance and local
source matching.

BGE-M3 supplies frozen text representations for semantic pattern
retrieval~\cite{chen2024bgem3}. It is neither a graph neural network nor a
jointly optimized component of the extractor. QLoRA enables compact adaptation
of the Qwen3 backbone~\cite{dettmers2023qlora,qwen2025qwen3}. Neither method
alone defines whether a candidate should enter the final assertion graph; the
proposed gates make that acceptance policy explicit.

\subsection{Evidence and safety-oriented knowledge construction}
Hallucination research distinguishes well-formed outputs from source-faithful
ones~\cite{ji2023hallucination}. In safety-related graph construction,
prompt-based processing of accident reports likewise illustrates the need to
trace generated knowledge back to its sources~\cite{liu2025llmaccidentkg}.
ADE offers a more controlled benchmark for a drug-safety-related relation
task~\cite{gurulingappa2012ade}. Nevertheless, an ADE label is an annotated
textual association, not a patient-level causal conclusion. The present study
evaluates extraction and implemented evidence constraints, rather than the
medical or operational validity of downstream decisions.

\section{Datasets and distribution analysis}
\label{sec:data}
\subsection{Benchmark sources and partitions}
We use the sentence-level ADE and CoNLL04 files distributed in SpERT format.
Counts refer to that preprocessed distribution rather than every sentence in
the original source corpora.
CoNLL04 retains its supplied training, development, and test partitions. The
ADE files contain 3,845 training sentences and 427 test sentences; a
deterministic development slice of 384 sentences is removed from the supplied
training pool, leaving 3,461 training sentences. The slice is selected by
sorting SHA-256 hashes of the string formed from the fixed split seed
\texttt{public-full-42} and each original record identifier. It is not selected
by validation performance. This is one fixed ADE split, not an average over
the multiple folds sometimes used in the literature.

Table~\ref{tab:data} gives counts recomputed from the exact imported artifacts.
The unit called a document in the software is one benchmark sentence, not a
whole medical case report or news article. Both datasets are English; the
results therefore support neither multilingual nor cross-lingual claims.
Published annotations are imported rather than locally re-annotated.

\begin{table}[htbp]
\centering\small\setlength{\tabcolsep}{4pt}
\caption{Dataset inventory for the exact experimental partitions. Length is
the number of whitespace-delimited tokens in the reconstructed source text.}
\label{tab:data}
\begin{tabular}{llrrrr}
\toprule
Dataset & Split & Sentences & Entities & Relations & Mean length \\
\midrule
ADE & train & 3461 & 8717 & 5470 & 21.17 \\
ADE & validation & 384 & 996 & 627 & 21.82 \\
ADE & test & 427 & 1126 & 724 & 21.24 \\
CoNLL04 & train & 922 & 3377 & 1283 & 28.77 \\
CoNLL04 & validation & 231 & 893 & 343 & 30.27 \\
CoNLL04 & test & 288 & 1079 & 422 & 28.94 \\
\bottomrule
\end{tabular}

\end{table}

\subsection{Entity and relation distributions}
ADE contains \texttt{Drug} and \texttt{Adverse-Effect} entities and a single
\texttt{Adverse-Effect} relation label. The imported endpoint convention is
\emph{adverse effect to drug}; we preserve it rather than reversing it to a
more intuitive drug-to-effect direction. CoNLL04 contains \texttt{Loc},
\texttt{Org}, \texttt{Peop}, and \texttt{Other} entities, with
\texttt{Work\_For}, \texttt{Kill}, \texttt{OrgBased\_In},
\texttt{Live\_In}, and \texttt{Located\_In} relations.

The training-label distributions are summarized in the accompanying tables, and
Table~\ref{tab:labels} lists exact label counts for every partition.
The single ADE relation label eliminates inter-class relation confusion but
does not eliminate endpoint detection, direction, or false-pair errors.
CoNLL04 introduces a broader relation inventory and more heterogeneous entity
roles. Consequently, a higher aggregate F1 on ADE should not be interpreted as
evidence that drug-safety reasoning is inherently easier than general-domain
reasoning; the supervised extraction tasks differ in structure and size.

\begin{table}[htbp]
\centering\small\setlength{\tabcolsep}{4pt}
\caption{Exact entity and relation annotation counts in the imported files.}
\label{tab:labels}
\begin{tabular}{lllrrr}
\toprule
Dataset & Kind & Label & Train & Validation & Test \\
\midrule
ADE & Entity & ADVERSE EFFECT & 4631 & 529 & 616 \\
ADE & Entity & DRUG & 4086 & 467 & 510 \\
ADE & Relation & ADVERSE EFFECT & 5470 & 627 & 724 \\
CoNLL04 & Entity & LOCATION & 1219 & 322 & 427 \\
CoNLL04 & Entity & ORGANIZATION & 616 & 170 & 198 \\
CoNLL04 & Entity & OTHER & 455 & 118 & 133 \\
CoNLL04 & Entity & PERSON & 1087 & 283 & 321 \\
CoNLL04 & Relation & KILL & 179 & 42 & 47 \\
CoNLL04 & Relation & LIVE IN & 330 & 91 & 100 \\
CoNLL04 & Relation & LOCATED IN & 247 & 65 & 94 \\
CoNLL04 & Relation & ORGANIZATION BASED IN & 271 & 76 & 105 \\
CoNLL04 & Relation & WORK FOR & 256 & 69 & 76 \\
\bottomrule
\end{tabular}

\end{table}

\subsection{Sentence lengths and coverage}
The mean training lengths are 21.17 reconstructed tokens for ADE and 28.77 for
CoNLL04. Their test means are 21.24 and 28.94, respectively. The cumulative
distributions describe all three partitions rather
than only the training set. These lengths are descriptive whitespace-token
counts, not Qwen tokenizer counts. The largest observed reconstructed length
is 97 for ADE and 118 for CoNLL04 across the imported partitions.

Every imported sentence contains at least one annotated relation. This is an
important property of the selected benchmark files: they do not measure
false-positive behavior over a deployment stream dominated by sentences with
no relevant relations. We retain the supplied benchmark composition and do
not create additional negative sentences for the reported tests.

\subsection{Split overlap and graph-context isolation}
The graph is built separately for each dataset from its training annotations
only. An exact-overlap check compares case-folded source strings after
collapsing whitespace. ADE has no train/development or train/test exact-text
overlap under this check. CoNLL04 has 14 development sentences and 19 test
sentences whose text matches training content. For these sentences, both
exact-anchor and high-recall graph context are cleared before generation.

Clearing graph hints does not undo exposure of the trained generator to
duplicate training text. We therefore retain the supplied partitions for the
main comparison and report an additional fixed-prediction test analysis that
excludes the 19 overlapping CoNLL04 test sentences. We do not claim that the
benchmark is duplicate-cluster isolated or that this exact-match check detects
all near duplicates.

\section{Proposed method}
\label{sec:method}
\subsection{Source-grounded representation}
Let a sentence $x$ be the deterministic single-space concatenation of its
released tokens. An entity mention is represented as
\begin{equation}
e=(b,u,y),\qquad x[b:u]=m,
\label{eq:entity}
\end{equation}
where $[b,u)$ is a character interval, $y$ is an entity label, and $m$ is the
copied mention text. A relation $r=(e_s,z,e_t)$ has typed directed endpoints
and relation label $z$. Source identifiers and evidence spans are stored with
these objects. Repeated strings at different offsets remain distinct mentions.
The character-span requirement in Equation~\eqref{eq:entity} prevents a
normalized name alone from receiving credit for the wrong occurrence.

The imported benchmark assigns every relation the status \texttt{explicit}
and uses its source sentence as the evidence span. This is a representation
convention, not a new annotation of minimal causal evidence or uncertainty.
Since both endpoints occur in the same input sentence, evidence co-occurrence
is a relatively weak condition in these experiments.

\subsection{Architecture and comparison systems}
Figure~\ref{fig:architecture} summarizes the implemented flow. SOE
(Source-Only Extraction), EAE (Exact-Anchor Extraction), and HRGE (High-Recall
Graph Extraction) are separately trained QLoRA adapters on the same Qwen3-4B
snapshot. SOE uses source text and the training-derived ontology without graph
hints. EAE receives conservative concept hints, while HRGE receives bounded
anchors, edge priors, and semantic relation patterns.

EVGE (Evidence-Verified Graph Extraction) applies the relation verifier to
HRGE. CFE (Conservative-Fusion Extraction) pairs accepted HRGE entities with
endpoint-filtered raw HRGE relations. PGE (Provenance-Gated Extraction) pairs
the same accepted entities with endpoint-filtered verified relations. These
three systems are deterministic post-processing variants, not additional
trained networks. In deployment PGE needs EAE and HRGE candidates; SOE is a
comparison branch and does not participate in fusion.

\begin{figure}[htbp]
\centering\includegraphics[width=\textwidth]{../figures/methodology_detailed_draft.pdf}
\caption{Detailed methodology for source-grounded Qwen3-4B extraction with
QLoRA adaptation, dataset-specific training graphs, deterministic evidence
gates, provenance-preserving graph output, and automatic audit records.}
\label{fig:architecture}
\end{figure}

\subsection{Provenance-controlled graph context}
Each training graph stores mention strings, types, directed relations, and
the sentence identifiers supporting them. For a training sentence $d$, a
candidate is excluded if its only supporting provenance is $d$. Candidates
with other supporting sentences remain eligible; their stored aggregate
support statistics are not recomputed as a fully leave-one-out graph. This
distinction prevents interpreting the provenance policy as complete removal
of every contribution from the current example.

EAE uses at most 20 exact-matching concept hints with type balancing. HRGE
uses three context channels. First, source-matched anchors must meet an
empirical type-purity threshold of 0.8 and an English string-length threshold
of four characters. Second, a typed edge prior requires support outside the
current example and both endpoint strings in the current source. Third,
frozen BGE-M3 retrieves semantically related relation-pattern records. If
$\mathbf{H}$ contains normalized eligible record embeddings and
$\mathbf{q}$ is the normalized source embedding, retrieval scores are
\begin{equation}
\mathbf{s}=\mathbf{H}\mathbf{q},\qquad
\mathcal I=\operatorname{TopK}_{k}\{i:s_i\geq\tau\}.
\label{eq:retrieval}
\end{equation}
Equation~\eqref{eq:retrieval} ranks candidate context, not accepted facts.
The frozen HRGE limits are 12 anchors with at most two per type, six edge
priors, and four semantic patterns with cosine similarity at least 0.72.
The prompt explicitly identifies these objects as hints. Exact-overlap
quarantine overrides retrieval and clears all graph context for the affected
evaluation sentences.

\subsection{Compact generation and deterministic reconstruction}
The generator emits entity identifiers, copied text and types, plus relation
identifiers, directed endpoint references, types, and explicit status. It is
not asked to calculate source offsets. Deterministic expansion assigns copied
mentions to occurrences in the supplied sentence and reconstructs character
spans; unmatched strings and invalid references are recorded. This separates
language-model candidate generation from source-coordinate bookkeeping.

All three generators use one training epoch, batch size one, gradient
accumulation four, and learning rate $2\times10^{-4}$. QLoRA uses rank eight,
alpha 16, dropout 0.05, NF4 double quantization, and bfloat16 computation.
Attention query/key/value/output projections and feed-forward gate/up/down
projections are adapted; the backbone remains frozen. AdamW uses the PyTorch
defaults for its other parameters, without a learning-rate scheduler or
warm-up. The maximum training sequence is 4,096 tokens, and prompt tokens are
masked from the loss. Training uses the same dataset partition for all three
branches and seed 42 throughout the primary comparisons.

Inference uses the recorded 4,096-token input limit, 1,024-token generation
budget, and 90-second job limit, with two inference workers on the local GPU.
The backbone revision and all input hashes are recorded in the artifact
manifest. Failed jobs are represented by empty predictions, not removed from
the denominator. These experiments use the full supplied sentence, without
annotation-guided rescue windows or label-dependent sentence truncation.

\subsection{Relation verification and entity acceptance}
For each relation label $z$, training annotations define permitted source-type
and target-type sets $S_z$ and $T_z$. The implemented type constraint is
$y_s\in S_z$ and $y_t\in T_z$; it does not additionally require that every
specific source/target type pair has been observed together. The verifier
$V(r)$ also requires existing distinct entity references, a known relation
label, explicit status, valid evidence, and both endpoint strings in a common
source-evidence span. Its accepted set is
\begin{equation}
R_V=\{r\in R_H:V(r)=1\}.
\label{eq:verifier}
\end{equation}
The entity gate uses three Boolean signals:
\begin{equation}
A(e)=A_{\rm agree}(e)\lor A_{\rm anchor}(e)\lor A_{\rm endpoint}(e),
\quad E_A=\{e\in E_H:A(e)=1\}.
\label{eq:accept}
\end{equation}
Agreement means that EAE predicts the same case-folded, whitespace-normalized
text and type. Anchor support requires a same-type source-gated HRGE anchor.
Endpoint support means that the entity participates in a relation retained by
Equation~\eqref{eq:verifier}. The three signals in Equation~\eqref{eq:accept}
are complementary but not statistically independent; verified endpoints are
a rule-based signal, not a second human judgement. Duplicate normalized
text/type candidates are collapsed within the sentence.

Let $F(R,E)$ retain relations in $R$ only when both referenced entities survive
in $E$. The two compositions are
\begin{equation}
G_{\rm CFE}=(E_A,F(R_H,E_A)),\qquad
G_{\rm PGE}=(E_A,F(R_V,E_A)).
\label{eq:composition}
\end{equation}
Equation~\eqref{eq:composition} makes the difference between raw-relation
fusion and verified fusion explicit. Accept/reject decisions are logged. A
retained relation satisfies the stated rules, but the rules alone cannot
establish whether the sentence asserts the proposed meaning.

\section{Experimental protocol}
\label{sec:experiment}
\subsection{Research questions and scope}
We ask whether PGE improves strict entity and relation extraction over SOE;
how graph context and deterministic selection contribute on development data;
and whether observed gains persist after exact-overlap sentences are removed.
External baselines establish whether the generative design is competitive
with other extraction families, rather than assuming that improvements over
SOE imply state-of-the-art performance.

The two-benchmark reporting scope was narrowed retrospectively after validation
and available test results had been inspected. It is a focused analysis, not
a preregistered or representative sample of all attempted extraction domains.
Within the existing protocol, model and gate settings were frozen
before test release. SOE and PGE were promoted when validation relation F1
did not regress by more than one percentage point and the required completion
and evidence checks passed. This was a pragmatic point-estimate promotion
rule, not a statistical non-inferiority test. The current paper uses the
completed seed-42 tests and does not infer multi-seed stability from them.

\subsection{Baselines and evaluation parity}
The internal development comparison includes all six named systems. The test
table includes the two promoted systems and three completed external
comparisons. Qwen3-4B zero-shot uses the same backbone without local adapters,
followed by the recorded source expansion and relation verifier. SpERT is
trained from \texttt{bert-base-cased} on the same training partitions for
20 epochs, with batch size two and a relation threshold of 0.4. Its final
checkpoint is used without selecting a test-best epoch. GLiNER+GLiREL uses
predicted, not gold, entities and relation calibration confined to an inner
training subset; the entity threshold is 0.5. These are heterogeneous baseline
configurations, not equal-compute experiments.

The recorded GLiREL calibration selects canonical relation labels and a
relation threshold of 0.0 for ADE, and naturalized labels with threshold 0.1
for CoNLL04. No development or test gold is used to choose those settings.

All reported main-table predictions are rescored with one public
character-span evaluator. An entity key contains its exact source interval
and type. A relation key contains both typed directed entity keys and the
relation label. Predicted evidence offsets must resolve to the supplied source;
the evaluator does not recover missing offsets by searching normalized text.
Duplicate keys receive one vote, invalid objects remain unmatched predictions,
and omitted jobs remain empty. For pooled counts,
\begin{equation}
P=\frac{TP}{TP+FP},\quad R=\frac{TP}{TP+FN},\quad
F_1=\frac{2TP}{2TP+FP+FN}.
\label{eq:metrics}
\end{equation}
Equation~\eqref{eq:metrics} is applied across sentences within each dataset,
not after averaging their individual F1 scores. Because all imported statuses
are explicit, claim-status-aware scores do not form a separate uncertainty
benchmark. The unified evaluator removes small denominator differences caused
by legacy duplicate handling; unchanged source artifacts and a reconciliation
report preserve the earlier metrics.

\subsection{Statistical and compliance analyses}
We compute 20,000 paired bootstrap resamples of test sentences, using seed
20260830 and fixed predictions, and recompute pooled F1 within each resample.
The percentile interval describes sentence-sampling uncertainty for the
completed model pair, not variability across training seeds. Intervals are
descriptive, unadjusted for multiple comparisons, and are not used to change
the frozen models or gates.

Evidence coverage and rule-invalid counts are measured separately from
extraction F1. These compliance counts refer to emitted prediction objects
before the evaluator's key deduplication. Exact-overlap sensitivity removes
the same affected test sentence identifiers from both systems, without
retraining or regenerating predictions. Test labels are used for reporting
statistics and scoring only, not for constructing graph hints or choosing
thresholds in this revision.

\section{Results}
\label{sec:results}
\subsection{Completed test comparisons}
Table~\ref{tab:test} reports the completed test results under the common
evaluator. PGE raises ADE relation precision from 80.62\% to 87.00\% and
recall from 67.82\% to 72.10\%, giving relation F1 of 78.85\% compared with
73.67\% for SOE. Entity F1 remains essentially unchanged, at 87.93\% versus
87.97\%. Thus, the main improvement is relational rather than a general gain
in mention recognition.

On CoNLL04, PGE raises relation precision from 54.61\% to 62.40\% and recall
from 56.16\% to 57.82\%. Relation F1 increases from 55.37\% to 60.02\%, while
entity F1 decreases from 83.83\% to 83.20\%. The composition trades some
entity coverage for a better relation precision--recall balance in this run.

\begin{table}[htbp]
\centering\scriptsize\setlength{\tabcolsep}{2.6pt}
\caption{Completed test comparisons, strict source-character-span
micro scores (\%). Qwen3-ZS is Qwen3-4B zero-shot with the recorded verifier;
GLiNER+GLiREL uses training-only calibration. Each dataset uses its own test
partition. Locally trained systems use seed 42; all rows are rescored with
the same evaluator.}
\label{tab:test}
\begin{tabular}{llrrrrrr}
\toprule
Dataset & System & Ent. P & Ent. R & Ent. F1 & Rel. P & Rel. R & Rel. F1 \\
\midrule
ADE & Qwen3-ZS + verifier & 63.83 & 63.32 & 63.58 & 39.36 & 22.24 & 28.42 \\
ADE & GLiNER + GLiREL & 79.30 & 64.65 & 71.23 & 60.31 & 43.23 & 50.36 \\
ADE & SOE & 92.21 & $\mathbf{84.10}\uparrow$ & $\mathbf{87.97}\uparrow$ & 80.62 & 67.82 & 73.67 \\
ADE & PGE & $\mathbf{92.89}\uparrow$ & 83.48 & 87.93 & $\mathbf{87.00}\uparrow$ & $\mathbf{72.10}\uparrow$ & $\mathbf{78.85}\uparrow$ \\
CoNLL04 & Qwen3-ZS + verifier & 48.78 & 53.94 & 51.23 & 30.92 & 19.19 & 23.68 \\
CoNLL04 & GLiNER + GLiREL & 60.78 & 14.37 & 23.24 & 13.79 & 3.79 & 5.95 \\
CoNLL04 & SOE & 85.83 & $\mathbf{81.93}\uparrow$ & $\mathbf{83.83}\uparrow$ & 54.61 & 56.16 & 55.37 \\
CoNLL04 & PGE & $\mathbf{86.04}\uparrow$ & 80.54 & 83.20 & $\mathbf{62.40}\uparrow$ & $\mathbf{57.82}\uparrow$ & $\mathbf{60.02}\uparrow$ \\
\bottomrule
\end{tabular}

\end{table}

SpERT has the highest entity and relation F1 among these completed systems:
91.37\%/84.08\% on ADE and 89.02\%/69.72\% on CoNLL04. PGE therefore improves
the matched generative baseline but does not surpass this trained span-based
model. PGE also exceeds the recorded zero-shot and GLiNER+GLiREL pipelines on
both extraction tasks. The latter observation is specific to their implemented
label mapping and calibration; it is not a claim that every configuration of
those model families would perform worse.

\subsection{Paired differences and uncertainty}
Using unrounded counts, PGE-minus-SOE relation F1 is 5.18 percentage points
on ADE, with a paired 95\% interval of [0.31, 10.29], and 4.65 points on
CoNLL04, with an interval of [0.04, 9.12]. The positive point estimates are
consistent with the test-table improvements, but the intervals are broad and
their lower bounds are close to zero. We interpret them as evidence about
these completed test predictions, not as a general guarantee across seeds.
Entity-F1 differences are $-0.04$ points on ADE, interval [$-1.75$, 1.62], and
$-0.63$ on CoNLL04, interval [$-2.38$, 1.04]; neither supports an entity-F1
improvement claim.

\subsection{Development-set component analysis}
Table~\ref{tab:ablation} separates candidate generation from deterministic
selection on development data. On ADE, HRGE has relation F1 of 76.29\%, above
SOE's 73.16\%. EVGE is identical to HRGE, and PGE is identical to CFE, showing
that the relation verifier removes no additional relations in these variants.
The entity gate reduces HRGE recall and gives PGE relation F1 of 75.91\%.
It is therefore incorrect to attribute all ADE gains to verification or to
claim that every stage independently improves F1.

On CoNLL04, HRGE relation F1 is 58.26\%, below SOE's 60.24\%. Verification
raises it to 59.60\%, and the final composition reaches 60.19\%. This is
effectively unchanged from SOE at the development point estimate, rather than
a clear development-set gain. Entity precision rises after gating, while
recall and F1 decrease. The distinction between development and test behavior
is retained rather than replacing the ablation with a test-selected variant.

\begin{table}[htbp]
\centering\scriptsize\setlength{\tabcolsep}{3pt}
\caption{Development-set ablation under the common strict-span evaluator
(micro \%, seed 42). These rows characterize the frozen architecture and are
not additional promoted test systems.}
\label{tab:ablation}
\begin{tabular}{llrrrrrr}
\toprule
Dataset & System & Ent. P & Ent. R & Ent. F1 & Rel. P & Rel. R & Rel. F1 \\
\midrule
ADE & SOE & 89.25 & 82.53 & 85.76 & 78.82 & 68.26 & 73.16 \\
ADE & EAE & 90.49 & 81.22 & 85.61 & 80.84 & 67.30 & 73.46 \\
ADE & HRGE & 90.35 & 85.54 & 87.88 & 80.04 & 72.89 & 76.29 \\
ADE & EVGE & 90.35 & 85.54 & 87.88 & 80.04 & 72.89 & 76.29 \\
ADE & CFE & 90.48 & 83.94 & 87.08 & 81.39 & 71.13 & 75.91 \\
ADE & PGE & 90.48 & 83.94 & 87.08 & 81.39 & 71.13 & 75.91 \\
CoNLL04 & SOE & 84.94 & 82.08 & 83.49 & 61.33 & 59.18 & 60.24 \\
CoNLL04 & EAE & 83.71 & 82.87 & 83.29 & 58.10 & 60.64 & 59.34 \\
CoNLL04 & HRGE & 82.35 & 81.52 & 81.94 & 60.06 & 56.56 & 58.26 \\
CoNLL04 & EVGE & 82.35 & 81.52 & 81.94 & 62.99 & 56.56 & 59.60 \\
CoNLL04 & CFE & 84.97 & 78.50 & 81.61 & 63.37 & 55.98 & 59.44 \\
CoNLL04 & PGE & 84.97 & 78.50 & 81.61 & 65.08 & 55.98 & 60.19 \\
\bottomrule
\end{tabular}

\end{table}

\subsection{Rule compliance and its interpretation}
The test SOE output contains 31 rule-invalid relation objects among 610 ADE
objects and 16 among 436 CoNLL04 objects. PGE retains 600 and 392 relation
objects, respectively, with zero rule-invalid objects. HRGE already has zero
invalid ADE relations and only three invalid CoNLL04 relations among 413
objects. These counts distinguish the contribution of candidate quality from
that of the verifier.

The measured unsupported-evidence count is zero for SOE, HRGE, and PGE on both
test datasets. Because each relation's evidence is the whole sentence,
co-occurrence is easy to satisfy once endpoints resolve. Thus, these
experiments do not demonstrate a reduction from a nonzero unsupported-claim
rate. They show explicit evidence binding and rule compliance, while the
gold-referenced F1 scores measure a different property: whether the predicted
typed relation agrees with the annotation.

\subsection{Exact-overlap sensitivity}
Removing the 19 exact train/test overlaps leaves 269 CoNLL04 test sentences,
1,023 gold entities, and 403 gold relations. SOE relation F1 becomes 53.79\%
and PGE becomes 59.10\%, a 5.31-point difference; entity F1 is 83.02\% and
82.90\%, respectively. The observed relation advantage therefore persists
without those exact overlaps. This analysis addresses exact repeated text,
not all forms of semantic overlap or pretraining contamination. ADE has no
such overlaps, so its sensitivity result is identical to the full test.

\section{Discussion}
\label{sec:discussion}
\subsection{What the two benchmarks establish}
The common result is a relation-F1 improvement over source-only generation,
not universal superiority over extraction architectures. The gains arise
alongside higher relation precision on both tests, while entity F1 stays
similar or declines slightly. This is consistent with an architecture designed
to govern candidate acceptance rather than simply maximize output volume.
The validation ablation further shows that graph conditioning and selection
interact: recall-oriented context can help on ADE but does not automatically
improve CoNLL04, and a verifier may be inactive when candidates already satisfy
its rules.

The distribution analysis helps interpret these differences. ADE has a larger
training pool, shorter average sentences, and a single relation class.
CoNLL04 has fewer training sentences, longer inputs, and five relation labels.
These factors are confounded rather than experimentally isolated. Their
association with the observed scores motivates separate dataset reporting;
it does not establish a causal explanation for the performance gap.

\subsection{Auditability as a distinct design objective}
A system can be less accurate than SpERT while still exposing a useful
acceptance mechanism for generated graph objects. The contribution is strongest
when the entire path is inspectable: graph support identifies why a candidate
was suggested; reconstructed offsets identify where it appears; gate records
identify why it was retained or rejected. However, this benefit is architectural.
No user study in the present experiments measures reduced analyst effort,
improved decisions, or greater trust. It would be inappropriate to substitute
perfect rule compliance for such evidence.

The gates also impose a coverage cost. Agreement operates on normalized
text/type pairs, and repeated mentions may be collapsed during fusion even
though the evaluator distinguishes offsets. This can remove legitimate
occurrences. A legal endpoint signature and sentence-level co-occurrence are
insufficient to prove a relation's meaning, and three gate signals should not
be interpreted as three independent confirmations. These observations explain
why the final output remains a candidate assertion graph for review.

\subsection{Implications for safety-related applications}
ADE links the experiments to drug-safety information extraction, while CoNLL04
tests general relational structure. The framework can present an analyst with
the candidate relation, source sentence, typed endpoints, and provenance of
its acceptance path. An analyst can then accept, amend, or reject the proposed
association. Absence from the retained graph is not evidence of the absence
of an adverse effect or another real-world relation.

The reported task stops before pharmacovigilance judgement, patient-specific
causal attribution, or operational action. Neither dataset validates risk
propagation, resilience assessment, or autonomous intervention. Extending this
work toward those applications requires task-specific outcome measures and
expert assessment in addition to extraction F1.

\section{Limitations and responsible use}
\label{sec:limitations}
The focused, retrospectively narrowed dataset scope and single completed
training seed limit generalization. The fixed ADE partition is not directly
interchangeable with published multi-fold results. Both corpora are English
and consist of relation-bearing sentences; multilingual behavior, long-report
processing, and false-positive rates on unrestricted text streams are outside
the evaluated setting. Sentence-level bootstrap intervals cannot recover
dependence between sentences originating from the same upstream article when
those article identifiers are unavailable.

Exact-overlap quarantine prevents direct graph hints for known repeated source
text but does not remove the generator's training exposure. The sensitivity
analysis excludes known exact duplicates only. Frozen pretrained encoders and
generators may also have seen related material during pretraining, which the
available artifacts cannot rule out. The ontology's allowed source and target
sets are less restrictive than a catalogue of observed joint type pairs.

The one-epoch generative systems, 20-epoch SpERT, and pretrained label-conditioned
models do not form a compute-matched benchmark. PGE requires two adapted
generators plus frozen retrieval, so compact adaptation must not be confused
with measured end-to-end cost savings. No new clinical annotations or
independent evidence-judgement labels were collected for this paper. Public
availability of a dataset or model does not by itself permit every downstream
redistribution or commercial use; original licenses remain applicable.

\section{Conclusions}
\label{sec:conclusion}
This study presents a provenance-preserving, evidence-gated graph-augmentation
framework and evaluates it on ADE and CoNLL04. Dataset statistics expose the
label and input-length differences underlying the two benchmarks. Under a
unified strict-span evaluator, the completed PGE tests improve relation F1
over SOE by 5.18 and 4.65 percentage points, respectively, while offering an
explicit source-linked acceptance path. The CoNLL04 relation advantage remains
after exact train/test overlaps are excluded. SpERT remains stronger in overall
extraction F1, and the development ablations show that the gates are not
uniformly beneficial in isolation. The supported outcome is controlled,
auditable graph-informed extraction for human review, rather than a claim of
universal accuracy dominance or validated autonomous safety reasoning.

\section*{Acknowledgements and funding}
Author-identifying funding and acknowledgement information is supplied on the
separate title page for double-blind review.
\section*{Declaration of competing interest}
The authors declare that they have no known competing financial interests or
personal relationships that could have appeared to influence this work.
\section*{Data and code availability}
The experiments use the public ADE and CoNLL04 distributions in SpERT format.
An anonymized reproducibility package will contain the split construction,
dataset statistics, source hashes, evaluation code, and derived result tables.
Dataset text and model weights remain subject to their original licenses.
The identifying repository URL is withheld during double-blind review.
\section*{Declaration of generative AI and AI-assisted technologies in the writing process}
The authors used OpenAI Codex to assist with code inspection, language editing,
figure preparation, and manuscript organization. The authors are responsible
for reviewing the resulting material and for the final content. The language
models evaluated as research systems are described separately in the Methods.

\bibliographystyle{elsarticle-num}
\bibliography{references}
\end{document}
